\origpage{288}{ROBERTSON AND SEYMOUR}
In particular, $\mathcal{D}(h, 0)$ is well-rooted. It follows from (3.3) that $\mathcal{C}(h, 0)$ is well-rooted; for if $(G, \emptyset) \in \mathcal{C}(h, 0)$ there are grafts $(G_1, \emptyset), \dots, (G_r, \emptyset) \in \mathcal{D}(h, 0)$ such that
\begin{equation*}
(G, \emptyset) = (G_1, \emptyset) \oplus \dots \oplus (G_r, \emptyset).
\end{equation*}
(In contrast, for $k > 0$ the class of all grafts expressible as the ``disjoint union'' of members of $\mathcal{D}(h, k)$ is not well-rooted, by (4.1).) Let $\mathcal{E}(h, k)$ be $\mathcal{D}(h, k) \oplus \dots \oplus \mathcal{D}(h, k)$, with $k$ summands. By (13.1), $\mathcal{E}(h, k)$ is well-rooted, and hence by (13.1) again, so is $\mathcal{E}(h, k) \oplus \mathcal{C}(h, 0)$.

Let $(G, T) \in \mathcal{C}(h, k)$, and let $\Gamma$ be a drawing of $(G, T)$ in some surface $\Sigma$ with $g(\Sigma) \le h$. Then $V(\Gamma) \cap bd(\Sigma') \ne \emptyset$ for at most $k$ components $\Sigma'$ of $\Sigma$, since $|T| \le k$, and so we may express $(G, T)$ as $(G_1, T_1) \oplus (G_2, \emptyset)$ where $(G_1, T_1) \in \mathcal{E}(h, k)$ and $(G_2, T_2) \in \mathcal{C}(h, 0)$. It follows that
\begin{equation*}
\mathcal{C}(h, k) \subseteq \mathcal{E}(h, k) \oplus \mathcal{C}(h, 0)
\end{equation*}
and so $\mathcal{C}(h, k)$ is well-rooted, as required. $\square$

\sectionhead{\textsc{References}}

1. D. Archdeacon, ``A Kuratowski Theorem for the Projective Plane,'' Thesis, The Ohio State University, 1980. 2. D. Archdeacon and P. Huneke, A Kuratowski theorem for non-orientable surfaces, \textit{J. Combin. Theory Ser. B} \textbf{46} (1989), 173--231. 3. H. Glover, P. Huneke, and C. S. Wang, 103 graphs that are irreducible for the projective plane, \textit{J. Combin. Theory Ser. B} \textbf{27} (1979), 332--370. 4. G. Higman, Ordering by divisibility in abstract algebras, \textit{Proc. London Math. Soc.} \textbf{2} (1952), 326--336. 5. J. Kruskal, Well-quasi-ordering, the tree theorem, and V\'azsonyi's conjecture, \textit{Trans. Amer. Math. Soc.} \textbf{95} (1960), 210--225. 6. K. Kuratowski, Sur le probl\`eme des courbes gauches en topologie, \textit{Fund. Math.} \textbf{15} (1930), 271--283. 7. N. Robertson and P. D. Seymour, Graph minors. III. Planar tree-width, \textit{J. Combin. Theory Ser. B} \textbf{36} (1984), 49--64. 8. N. Robertson and P. D. Seymour, Graph minors. IV. Tree-width and well-quasi-ordering, \textit{J. Combin. Theory Ser. B} \textbf{48} (1990), 227--254. 9. N. Robertson and P. D. Seymour, Graph minors. V. Excluding a planar graph, \textit{J. Combin. Theory Ser. B} \textbf{41} (1986), 92--114. 10. N. Robertson and P. D. Seymour, Graph minors. VI. Disjoint paths across a disc, \textit{J. Combin. Theory Ser. B} \textbf{41} (1986), 115--138. 11. N. Robertson and P. D. Seymour, Graph minors. VII. Disjoint paths on a surface, \textit{J. Combin. Theory Ser. B} \textbf{45} (1988), 212--254.
